% ============================================================================
% 相关工作 - 02_related_work.tex
% ============================================================================

\section{相关工作}
\label{sec:related_work}

本节系统回顾无人机路径规划与风险评估领域的研究进展，从静态风险模型、动态环境建模以及多目标优化算法三个维度梳理现有工作，并指出现有研究的局限性。

\subsection{无人机第三方风险评估模型}


\subsection{动态环境因素建模}



\subsection{风险感知路径规划算法}


\subsection{研究缺口与本文贡献}

综上所述，现有研究存在以下关键局限：首先，多数研究采用常数坠机概率的静态风险假设，未能反映气象和建筑环境的动态影响；其次，现有模型往往仅关注物理撞击风险的单一维度，忽视噪声滋扰等社会影响成本；第三，人口和交通密度的时间潮汐效应未被充分纳入风险评估框架，时空耦合不足；最后，高维动态风险场中的路径搜索效率低下，算法可扩展性差，难以满足实时规划需求。

针对上述问题，本文提出一种多维时空耦合的第三方风险综合成本评估模型，并设计分层混合路径寻优策略。具体贡献包括：构建动态坠机概率模型，整合风场、降雨和城市建筑三重环境因子；开发时空潮汐引力模型，模拟动态人口和车辆暴露密度；建立多维度风险成本函数，涵盖致命风险、财产损失和噪声滋扰；提出EDA-K-means-CostA*分层算法框架，实现高效的全局-局部协同优化。
